\section{Discussion}
\label{sec:discussion}

\subsection{Relation to Beukers--Vlasenko's Conjecture 7.5}
\label{ssec:BV}

The natural home for a congruence like Conjecture~\ref{conj:master} is the theory of
Dwork crystals, and it is worth being precise about how the two fit together. We read
\cite{BVthree} from its source; its \S7 treats \emph{completely symmetric Calabi--Yau
families}, and its numbered items there are 7.1 (definition), 7.2 (lemma), 7.3 (theorem),
7.4 (remark) and 7.5 (conjecture).

The setting is $f(x)=1-tg(x)$ and $F(t)=\sum_nf_nt^n$ with $f_n=\ct\bigl(g(x)^n\bigr)$;
for the sporadic families these $f_n$ are exactly our $A(n)$. Writing $F_N(t)$ for the
truncation of $F$ at $t^N$, and $G(t)$ for the series appearing in the second solution
$F(t)\log t+G(t)$ of the (second-order) Picard--Fuchs operator, \emph{the coefficients of
$G$ are our $B(n)$}, up to normalisation. What \cite{BVthree} proves --- their eq.~(7.6),
from \cite[eq.~(7)]{BVtwo} --- is that for all $m,s\ge1$,
\begin{equation}\label{eq:BV76}
   \frac{F(t)}{F(t^\sigma)}\;\equiv\;\frac{F_{mp^s}(t)}{F_{mp^{s-1}}(t^\sigma)}\pmod{p^s}.
\end{equation}
Reading off coefficients with the naive lift $t^\sigma=t^p$ and $s=1$ gives the
Dwork--Mellit--Vlasenko congruence whose sporadic specialisation is precisely the
first-row Lucas congruence $A(ap+r)\equiv A(a)A(r)\pmod p$ of
Theorem~\ref{thm:arow}. Their Conjecture~7.5 asserts that \eqref{eq:BV76} holds modulo
$p^{2s}$ --- for $m=2$ in the hypercubic and hyperoctahedral families and $m=n+1$ in the
simplicial ones --- where $\sigma$ is the \emph{excellent} Frobenius lift, the unique lift
with $q^\sigma=\gamma^{p-1}q^p$ for the canonical (mirror) coordinate
$q(t)=t\exp\bigl(G(t)/F(t)\bigr)$; by their Theorem~7.3 that lift is characterised by the
vanishing of the off-diagonal Frobenius entry. The authors state that they could not prove
7.5 for a single $g$.

Four points relate this to Conjecture~\ref{conj:master}.

\begin{enumerate}[leftmargin=1.7em]
\item \emph{They are the two rows of one $2\times2$ Frobenius matrix.} The crystal
  $\mathrm{CY}(g)(2)$ of \cite{BVthree} is free of rank two on $1/f$, $\theta(1/f)$, and
  its Frobenius has diagonal $(\lambda_0,\lambda_2)$ with $\lambda_0=F(t)/F(t^\sigma)$ the
  unit root. Conjecture 7.5 is a statement about $\lambda_0$ alone, sharpened from modulo
  $p^s$ to modulo $p^{2s}$. Conjecture~\ref{conj:master} is the statement about the
  \emph{other} diagonal entry: that Frobenius acts on the second graded piece with
  eigenvalue $\chi(p)p^{w}$. In coefficient form, their $\lambda_0$-congruence \emph{is}
  the first row; ours \emph{is} the second row. Neither formally implies the other.
\item \emph{But 7.5's standing hypothesis is a weak form of our conclusion.} The excellent
  lift exists as a lift of $\Z_p[[t]]$ precisely because
  $q(t)=t\exp(G/F)\in\Z_p[[t]]$ --- $p$-integrality of the mirror map, which is a
  statement about $G$, i.e.\ about $B(n)$. Our Corollary~\ref{cor:integrality},
  $d_n^wB_n\in\Zp$, obtained by iterating the descent over base-$p$ digits, is exactly the
  coefficient-level shadow of that integrality, and it is sharper: it does not merely
  assert integrality of the mirror map, it pins the depth at exactly $w$ and exhibits the
  eigenvalue $\chi(p)$. So the implication runs
  \[
   \begin{gathered}
    \text{Conjecture~\ref{conj:master}}\;\Longrightarrow\;
    \text{(coefficient form of mirror-map integrality)}\\
    \Longrightarrow\;
    \text{(the excellent lift is defined over $\Z_p[[t]]$)},
   \end{gathered}
  \]
  which is the standing hypothesis of Conjecture 7.5. In this precise sense our conjecture
  is \emph{upstream} of theirs, and not equivalent to it.
\item \emph{The character is invisible in the unit-root formulation and must be added.}
  Beukers and Vlasenko work with $\lambda_0=F/F^\sigma$, where no character can appear:
  the unit root of the holomorphic period is a unit, and the first-row Lucas congruence is
  untwisted for all fifteen families --- \Verified\ here and proved in
  \cite{MalikStraub}. The character surfaces only on the \emph{non-unit} eigenvalue. Any
  crystal-theoretic proof of Conjecture~\ref{conj:master} must therefore carry the
  information that, for seven of the fifteen families, the second Frobenius eigenvalue is
  $\chi(p)p^{w}$ and not $p^{w}$ --- i.e.\ that the family is attached to a quadratically
  twisted (CM) modular parametrisation. This is a concrete correction to make to any
  slogan of the form ``Dwork crystals imply Lucas congruences for the second solution''.
\item \emph{What our conjecture does not give.} Nothing towards the modulo-$p^{2s}$
  doubling that is the actual content of 7.5. Our floor is exactly $w$ and provably not
  more in the sense that equality is attained in every cell tested, so there is no hidden
  supercongruence in the second row of the kind 7.5 predicts for the first.
\end{enumerate}

We note for completeness that we could not find any second-solution version of these
congruences stated in \cite{BVthree} or in Vargas-Montoya's work
\cite{VMcanon,VMmum}; \cite{VMmum} proves Lucas congruences explicitly and exclusively for
the holomorphic solution, and \cite{VMcanon} proves $p$-integrality of $\exp(y_1/y_0)$
with no congruence attached.

\subsection{Frobenius, not the archimedean limit}
\label{ssec:frobenius}

Observation~\ref{obs:nolimit} deserves emphasis. Three sporadic sequences --- $\mathbf B$,
$\delta$, $\eta$ --- have characteristic roots forming a complex-conjugate pair of equal
modulus, so Poincar\'e--Perron gives no dominant solution and $B(n)/A(n)$ oscillates
forever: there is no Ap\'ery limit to speak of. Nonetheless all three satisfy
Conjecture~\ref{conj:master}, with a flat floor, and with characters $\chi_{-3}$, trivial
and $\chi_5$ respectively.

The $p$-adic Ap\'ery limit therefore exists where the archimedean one does not. This is
the sharpest indication available to us that the descent is a statement about a Frobenius
eigenvalue and not about an asymptotic ratio: the archimedean limit is an accident of the
relative sizes of the characteristic roots, whereas the $p$-adic descent is insensitive to
them. It also means that ``the character of the Ap\'ery limit'' --- which is how we
identified $\chi$ in twelve of the fifteen cases --- cannot be the definition of $\chi$;
for those three it must be read off from the descent itself, or from the modular
parametrisation.

\subsection{Beyond rank two}
\label{ssec:rank3}

The natural next target is the Brown--Zudilin cellular pair for $\zeta(5)$
\cite{BrownZudilin}, where the integral row is a double binomial sum $Q_n$, the second and
third solutions $P_n$ and $\widehat P_n$ satisfy the same order-three recurrence, and the
expected descent has $w=5$. The obstructions are structural and were identified early in
this programme: the local system has rank three, so the descent must be proved for a
three-step filtration at once, with the intermediate weight-three object $\widehat P_n$
appearing in the $P_n$ descent as a cross-term with no $\zeta(3)$ analogue; the summand is
not a perfect square, so the per-carry valuation slack of Remark~\ref{rem:square}
disappears; and the pole to be assembled is of order five rather than three, which roughly
squares the case analysis of Lemma~\ref{lem:V}. The first-row Lucas congruence
$Q_{ap+r}\equiv Q_aQ_r\pmod p$ is proved and formalised (Section~\ref{sec:lean}).

The lesson of Section~\ref{sec:minimal} applies here with force: the right search is not
``which weight-five harmonic monomials span $P_n$'', but ``which weight has an
$n$-difference equal to a single hypergeometric term times a rational function'' --- a
smaller, linear and decidable search, and one insensitive to the large kernel that makes
monomial fitting degenerate. And the top-cell constant, which at weight three is the
classical $5$ of Lemma~\ref{lem:D5}, is computable in advance from the residue at the top
of the support, giving a necessary numerical fingerprint that any candidate weight must
match.

\subsection{Open problems}
\label{ssec:open}

\begin{problem}\label{prob:master}
Prove Conjecture~\ref{conj:master}, in any of its three independent directions: the
multi-digit statement, the sharpening from modulo $p$ to modulo $p^w$, and the removal of
the tameness hypothesis. Even the weight-three case, Conjecture~\ref{conj:master3}, is
open in all three.
\end{problem}

\begin{problem}\label{prob:eps}
Prove Theorem~\ref{thm:brow} by an $\varepsilon$-deformation. Gessel's companion
congruence \cite[Thm.~1.4]{Beukers2211} is structurally the first parameter-derivative of
the first-row congruence; our $b_n$ is the weight-three companion, i.e.\ the third
derivative. Writing $a_n=\ct[\Lambda_\varepsilon^n]$ for a deformed Laurent polynomial,
applying the Dwork--Mellit--Vlasenko ratio congruence \cite{MellitVlasenko} to the
deformed family and taking $\partial_\varepsilon^3|_{\varepsilon=0}$ should land the
weight-three pole and multiply the modulus by $p^3$, since each filtration level costs one
power of $p$ \cite[Prop.~7.7]{BVthree}.
\end{problem}

\begin{problem}\label{prob:lemmaD}
Supply the valuation bound (``a Kummer carry in a wide binomial beats the order-one pole
of a harmonic letter at $\lfloor2k/p\rfloor$'') that would remove the tameness hypothesis
(H4) and put $\alpha$, $\varepsilon$, $s_7$ and $\mathbf E$ inside Theorem~\ref{thm:LB}.
This would also produce the first \emph{proved} instance of the twisted case $e=1$, which
at present has no tame example.
\end{problem}

\begin{problem}\label{prob:seven}
Find harmonic decompositions for the remaining seven families. For $\mathbf C$ this
provably requires letters that its own summand does not supply
(Observation~\ref{obs:negatives}); a constant-term representation \cite{Gorodetsky} is
the natural route. For $\delta$, letters attached to the $3$-section of the summand.
\end{problem}

\begin{problem}\label{prob:zeta2}
Close the $\zeta(2)$-Ap\'ery case of Section~\ref{ssec:obstruction}: either exhibit a
gauge (a shift by an element of the $11$-dimensional kernel of the fitting system) in
which the correction telescopes, or run the argument in a $\Pi\Sigma$ field with nested
sums. Proposition~\ref{prop:obstruction} shows no rational $\times$ harmonic-monomial
ansatz can work.
\end{problem}

\begin{problem}\label{prob:chi-intrinsic}
Give an intrinsic definition of $\chi$ --- one that does not go through the Ap\'ery limit,
and so applies to $\mathbf B$, $\delta$, $\eta$ --- and prove that it is the character of
the quadratic twist of the modular parametrisation attached to the family.
\end{problem}
